\section{What the Walk Is Actually Asking}
\label{sec:ch15-what-the-walk-is-asking}

\index{resolvability!the question behind the message|(}%
Chapter~\ref{ch:composition} described this check as a walk out from each root
field, and said the message is a trace of where the walk got stuck. That is
still true and it is not enough to read one of these, because it leaves the
impression that the printed path is the subject. The path is a byproduct. The
subject is a field, and the question is whether any route to it exists.

\index{resolvability!an existence question has one answer}%
An existence question has one answer, so it gets one report. Six broken routes
to one field produce one error about that field, carrying whichever route the
walk was on when it ran out of moves. Three fields, three errors. That is the
whole of why
section~\ref{sec:ch15-the-route-it-named-is-one-of-six} happens, and it is
worth holding onto, because the shape of the message argues against it: a query
with an arrow in it looks like a claim about that query.

\index{resolvability!the check stops at the first failing root field}%
The composer also stops. It walks the root fields in order and returns on the
first one that fails, so the rest of the schema is not examined in that run.
Fix what it names, compose again, and a second failure you have not seen yet
can be waiting behind the first. That is a property of this composer rather
than of the idea: the check as the Composite Schemas draft writes it is a stack
of paths, popped one at a time~\autocite{compositeschemasspec}, and an
algorithm shaped like that has no particular reason to keep going after an
answer. What each implementation does about it is the implementation's.

\subsection{What an Entity Ancestor Is}

\index{entity ancestor!the term the message uses}%
Two of the four reason lines use a phrase that appears nowhere in a schema:

\begin{minted}[fontsize=\footnotesize,breakanywhere]{text}
 - The entity ancestor "Session" in subgraph "sessions" has no accessible target entities (resolvable @key directives)
in the subgraphs where "Session.averageScore" is defined.
\end{minted}

\index{entity ancestor!is the last entity the walk stood on}%
An entity ancestor is the last entity the walk was standing on before it
reached the field it could not get to. Getting from there to the field means
crossing into another service, and the only crossing federation offers is a key
you are allowed to enter through. That is what an \emph{accessible target
entity} is: a copy of the same type, in a subgraph that declares a key without
\code{resolvable: false} on it. The fourth line then reports that no other
entity above this one offers a crossing either, which is the composer saying it
looked for a way round and found none.

\index{entity ancestor!need not be on the path it printed}%
What the phrase does not tell you is that the subgraph it names need not be on
the path above it. In the four-service error the path runs through
\code{search} and the ancestor is in \code{sessions}. Both services declare a
\code{Session}, both of them fail to reach ratings from it, and both are
correct. The name in the sentence is whichever of the two the composer visited
first, which is subgraph order in \code{graph.yaml}, which is why
section~\ref{sec:ch15-the-route-it-named-is-one-of-six} could change it by
reversing four lines of configuration.

\index{entity ancestor!read it as a type, not as an address}%
So read that line as a claim about the \emph{type} and treat the subgraph name
as a hint. \code{Session} cannot get to ratings from anywhere. Which of its
copies the composer was holding when it said so is not information about your
graph.

\subsection{Two Passes, and Only One of Them Is Load-Bearing}

\index{resolvability!still separable from merging}%
Chapter~\ref{ch:composition} established that composition is two passes and
that the second can be switched off, which is still the case: put
\code{--disable-resolvability-validation} on the command and all three errors
go away and a config is written. The command's help still carries the flag on
\code{wgc} 0.129.9, though the reference page for the command no longer lists
it and documents it under a different command instead.

\index{resolvability!a graph composed with the flag}%
A graph composed that way is a graph where some field answers at run time with
something no composition warned about, and
section~\ref{sec:ch15-three-fields-three-failures} is a catalog of what that
looks like. I would not use the flag. What matters more is the next section,
where the check reaches the same outcome on its own.
\index{resolvability!the question behind the message|)}%
